\documentclass{article}
\usepackage{amsmath,amssymb}
\usepackage[vocab={realsets}]{leanify}
\newtheorem{lemma}{Lemma}
\begin{document}

\section{Notation}

We write $\overline{S}$ for the closure of $S$.
We write $S^\circ$ for the interior of $S$.
We write $\lambda(S)$ for the Lebesgue measure of $S$.
We write $\delta(S)$ for the diameter of $S$.
We write $\sigma(S)$ for the supremum of $S$.
We write $\mathcal{C}$ for the Cantor set.
We write $Q$ for the set of rational numbers.

\section{Sets of real numbers}

\begin{lemma}\label{heine}
Let $S$ be a set. Then $S$ is compact if and only if $S$ is closed and bounded.
\end{lemma}

\begin{lemma}\label{closure}
For every set $S$, $\overline{S}$ is closed and $S^\circ$ is open.
\end{lemma}

\begin{lemma}\label{supremum}
Let $S$ be a nonempty compact set. Then $\sigma(S)\in S$.
\end{lemma}

\begin{lemma}\label{rationals}
$Q$ is dense and countable, and $Q$ is null.
\end{lemma}

\begin{lemma}\label{cantor}
$\mathcal{C}$ is compact, perfect, nowhere dense and uncountable. Moreover, $\lambda(\mathcal{C})=0$.
\end{lemma}

\begin{lemma}\label{countable}
Let $S$ be a countable set. Then $S$ is a null set.
\end{lemma}

\begin{lemma}\label{open}
Let $U$ be a nonempty open subset of $\mathbb{R}$. Then $U$ is uncountable, and there is $x\in U$ such that $x$ is rational.
\end{lemma}

\begin{lemma}\label{measure}
Let $S$ be a nonempty bounded measurable set. Then $\lambda(S)\le\delta(S)$.
\end{lemma}

\begin{lemma}\label{zero}
Let $S$ be a set of measure zero and let $T$ be the closure of $S$. Then $S^\circ$ is empty, and $S$ is of measure zero or $T$ is uncountable.
\end{lemma}

\begin{lemma}\label{sup-exists}
Let $A$ be a nonempty bounded subset of $\mathbb{R}$. Then $\sup A$ exists.
\end{lemma}

\begin{lemma}\label{sup-upper-bound}
Let $A$ be a nonempty set of real numbers that is bounded above. Then $\sup A$ is an upper bound of $A$.
\end{lemma}

\begin{lemma}\label{interval-compact}
The interval $[0,1]$ is compact.
\end{lemma}

\begin{lemma}\label{finite-maximum}
Let $A\subseteq\mathbb{R}$ be finite and nonempty. Then $A$ has a maximum.
\end{lemma}

\begin{lemma}\label{heine-borel}
Let $K$ be a compact subset of $\mathbb{R}$. Then $K$ is closed and bounded.
\end{lemma}

\begin{lemma}\label{non-empty-compact}
Let $K$ be a compact non-empty subset of $\mathbb{R}$. Then $K$ is bounded.
\end{lemma}

\begin{lemma}\label{set-elements}
Let $S$ be a finite set with $n$ elements. Then $|S|=n$.
\end{lemma}

\begin{lemma}\label{intervals}
Let $I$ be an open interval. Then $I$ is open. Every closed bounded interval is compact. Let $J$ be an interval. Then $J$ is convex.
\end{lemma}

\begin{lemma}\label{infimum}
Let $B$ be a nonempty bounded set of real numbers. Then $\inf_{z\in B}z\le\sup_{z\in B}z$.
\end{lemma}

\begin{lemma}\label{finite-sets}
Every nonempty set of natural numbers has a least element. Let $S$ be a nonempty set of integers that is bounded above. Then $S$ has a greatest element. Let $A$ and $B$ be finite sets. Then $|A\cup B|=|A|+|B|-|A\cap B|$. For all sets $A$, $B$ and $C$, $A\cap(B\cup C)=(A\cap B)\cup(A\cap C)$. Every countable subset of $\mathbb{R}$ has Lebesgue measure zero. Let $D$ be a countable set. Then the Lebesgue measure of $D$ is zero.
\end{lemma}

\begin{lemma}\label{inclusions}
Let $A\subseteq B$ be finite sets. Then $|A|\le|B|$. Let $C$ be a set and let $x$ be a real number. If $C$ contains $x$, then $C$ is nonempty.
\end{lemma}

\begin{lemma}\label{sequences-of-sets}
Let $(A_n)$ be a sequence of countable sets of real numbers. Then $\bigcup_{n=1}^{\infty}A_n$ is countable. Let $(B_n)$ be a sequence of sets of real numbers with $B_{n+1}\subseteq B_n$ for all $n$. Then $\bigcap_{n=0}^{\infty}B_n\subseteq B_0$ and $\bigcup_{n=0}^{3}B_n=B_0$.
\end{lemma}

\begin{lemma}\label{relative-clauses}
Every nonempty set of real numbers that is bounded above has a supremum. A nonempty set of integers that is bounded above has a greatest element.
\end{lemma}

\begin{lemma}\label{heine-borel}
A subset of $\mathbb{R}$ is compact if and only if it is closed and bounded.
\end{lemma}

\begin{lemma}\label{completeness}
Every nonempty subset of $\mathbb{R}$ that is bounded above has a least upper bound. Every nonempty set of real numbers that is bounded below has a greatest lower bound. Between any two distinct real numbers there is a rational number. There is no smallest positive real number.
\end{lemma}

\begin{lemma}\label{containing-sets}
Let $A$ be a set of real numbers containing $0$. Then $A$ contains $0$ and $A$ is nonempty. Let $B$ and $C$ be sets of real numbers with $B\subseteq C$. Then $C$ contains $B$. Then $(0,1)$ does not contain $0$.
\end{lemma}

\begin{lemma}\label{two-sets}
The union of two bounded sets of real numbers is bounded. The intersection of two closed sets of real numbers is closed. The union of two countable sets of real numbers is countable.
\end{lemma}

\end{document}
